\section*{第 \ref{ch:g6:exploring-our-environment} 章——探索我们的环境}

\begin{solution}{exo:g6:exploring-our-environment:1}
\omterm{def:g6:exploring-our-environment:components}{生物}（\omterm{prop:g4:plants-without-seeds:damp}{苔藓}、潮虫）；\omterm{def:g6:exploring-our-environment:components}{矿物质}（墙上的
石头、水坑）；\omterm{def:g6:exploring-our-environment:components}{生物}的痕迹和遗骸（枯叶、
蜗牛壳）。
\end{solution}

\begin{solution}{exo:g6:exploring-our-environment:2}
\omterm{def:g6:exploring-our-environment:conditions}{温度}（用\omterm{def:g6:exploring-our-environment:conditions}{温度}计测量，单位为摄氏度）；
空气和土壤的\omterm{def:g6:exploring-our-environment:conditions}{湿度}；光照，从全日照到深阴。
\end{solution}

\begin{solution}{exo:g6:exploring-our-environment:3}
\omterm{def:g6:exploring-our-environment:components}{生物}：\omterm{prop:g4:plants-without-seeds:damp}{蕨类}、蜈蚣。\omterm{def:g6:exploring-our-environment:components}{矿物质}：水坑、
墙上的石头。痕迹：蜗牛壳、鸟粪。
\end{solution}

\begin{solution}{exo:g6:exploring-our-environment:4}
相差 16 摄氏度（\qty{38}{\degreeCelsius}对
\qty{22}{\degreeCelsius}），以及全日照对深阴
（干对湿也是如此）。耐旱耐晒的杂草
与南墙根相匹配；\omterm{prop:g4:plants-without-seeds:damp}{苔藓}、\omterm{prop:g4:plants-without-seeds:damp}{蕨类}和潮虫
则与凉爽潮湿的墙角相匹配。
\end{solution}

\begin{solution}{exo:g6:exploring-our-environment:5}
它的外壳让水分散失，因此在干燥的阳光下，
它的身体会以致命的速度失去水分。这就是为什么
潮虫常见于潮湿、阴暗的\omterm{def:g6:exploring-our-environment:microhabitat}{微生境}——石头下、枯木下。
\end{solution}

\begin{solution}{exo:g6:exploring-our-environment:6}
例如：被啃过的榛子——林姬鼠；蜘蛛网——
蜘蛛本身；泥地上的足迹——苍鹭（鸟粪、\omterm{def:g1:animals-around-us:covering}{羽毛}、
壳也有同样作用）。
\end{solution}

\begin{solution}{exo:g6:exploring-our-environment:7}
在一个封闭的箱子里，一半铺湿纸、一半铺干纸，
把潮虫放在中线上；几分钟后，几乎全部
聚集在湿纸一侧（光与阴的情况同样如此）。
这说明每只\omterm{def:g1:animals-around-us:animal}{动物}都会停留在\omterm{def:g6:exploring-our-environment:conditions}{条件}适合
其身体的地方——整个类群的分布图无需人为驱赶即可自行形成。
\end{solution}

\begin{solution}{exo:g6:exploring-our-environment:8}
因为只有用相同方法测得的数据才可比较：
差异必须来自不同的地点，而非测量方法。
这是野外版的公平测试——只改变一个因素（地点），
其余\omterm{def:g6:exploring-our-environment:conditions}{条件}保持不变。
\end{solution}

\begin{solution}{exo:g6:exploring-our-environment:9}
石头顶面：有光照、干燥、温暖——地衣
（另有一只晒太阳的苍蝇作为访客）。石头底面：
阴暗、潮湿、凉爽——潮虫（或蜈蚣）。
两处相隔 5 厘米，却分别住着两组相互匹配的居民。
\end{solution}

\begin{solution}{exo:g6:exploring-our-environment:10}
第 1 步：划定一个确定的一平方米样地。第 2 步：测量——
记下\omterm{def:g6:exploring-our-environment:conditions}{温度}计读数，记录土壤\omterm{def:g6:exploring-our-environment:conditions}{湿度}（用手感受），
并记下光照等级。第 3 步：认真搜寻，翻开石头后
放回原处，三类成分全部列出。第
4 步：当场记录——记录数量，而非“一些”。第 5 步：
只有一个地点时根本无法比较；应以同样方法
再调查第二个地点。
\end{solution}

\begin{solution}{exo:g6:exploring-our-environment:11}
解释：常春藤能在北墙的\omterm{def:g6:exploring-our-environment:conditions}{条件}范围内生长（阴凉、
潮湿），却不能在南墙（强光、干旱）处生长。验证：
测量两面墙根的\omterm{def:g6:exploring-our-environment:conditions}{温度}、\omterm{def:g6:exploring-our-environment:conditions}{湿度}和光照，
并核对常春藤已知的偏好；
或在其他墙面上寻找常春藤，记录这些墙面
共同的朝向。
\end{solution}

\begin{solution}{exo:g6:exploring-our-environment:12}
为消除墙面差异，改在圆形箱中进行选择实验
（各处墙面相同），每次实验交换湿、干两半的位置——
第 1 次湿润在东侧，第 2 次湿润在西侧。
若潮虫总是跟随湿润的一侧，无论湿润在哪里，
便可排除墙面的影响；两次实验之间只改变了
被测\omterm{def:g6:exploring-our-environment:conditions}{条件}，其余\omterm{def:g6:exploring-our-environment:conditions}{条件}保持相同。
\end{solution}

\begin{solution}{pb:g6:exploring-our-environment:1}
\textbf{1.} 高大的树篱遮住了北侧沟渠（并为其遮挡），
而南侧沟渠整天面向阳光。

\textbf{2.} \omterm{def:g6:exploring-our-environment:conditions}{温度}：\qty{19}{\degreeCelsius}对
\qty{31}{\degreeCelsius}。\omterm{def:g6:exploring-our-environment:conditions}{湿度}：北侧沟渠的土壤潮湿（能粘
成团），南侧沟渠的土壤干燥（扬尘）。光照构成第三项差异：
阴凉对全日照。

\textbf{3.} \omterm{def:g6:exploring-our-environment:conditions}{条件}会随一天中的时间变化：到下午 6 点，
任一地点都会变凉，光照也会减弱，因此所测得的差异
会把地点与时间混在一起。在同一时间测量，
比较才公平。

\textbf{4.} 为使搜寻投入相等：若在某一地点投入更大面积或更多时间，
就会在那里发现更多\omterm{def:g6:exploring-our-environment:components}{生物}——
这样比较哪条沟渠更富饶便不公平。

\textbf{5.} 两份清单都只列出了\omterm{def:g6:exploring-our-environment:components}{生物}。缺少的是
\omterm{def:g6:exploring-our-environment:components}{矿物质}成分（土壤、石头、水）和痕迹
——这些也应当记入完整的清单中。

\textbf{6.} 潮虫：沟渠 N——它们需要潮湿和
遮荫，否则会干枯。蟾蜍：沟渠 N——\omterm{def:g1:the-five-senses:sense}{皮肤}湿润，
需要潮湿的庇护所。蜥蜴：沟渠 S——它靠晒太阳
取暖，并在能晒太阳的地方捕猎。

\textbf{7.} 蜗牛或许躲藏在那里生活（可能只在潮湿的夜间活动），
或许这只壳是来自别处或别的季节的遗骸——
痕迹只能证明曾在
\emph{某个时候}存在过，而不能说明现在仍在此居住。

\textbf{8.} 黑暗、潮湿、凉爽。依据
\cref{prop:g6:exploring-our-environment:decide}的规律：每个地点
都居住着其分布范围与该地\omterm{def:g6:exploring-our-environment:conditions}{条件}
相符的\omterm{def:g5:biodiversity:species}{物种}。

\textbf{9.} 沟渠 N 有\omterm{prop:g4:plants-without-seeds:damp}{苔藓}和蟾蜍，因为其遮荫使其
保持凉爽潮湿，这正是它们身体所需要的；沟渠 S 有蜥蜴和
叶片覆有蜡质的杂草，因为其充足的阳光使其炎热干燥，
这适合靠晒太阳取暖的捕食者和耐旱的\omterm{def:g1:plants-around-us:plant}{植物}。

\textbf{10.} 若没有树篱，沟渠 N 将受到充足的阳光照射：更热、
更干燥——\omterm{def:g6:exploring-our-environment:conditions}{条件}将逐渐接近今日的沟渠 S。
可以预料，喜湿的\omterm{def:g6:exploring-our-environment:components}{生物}（\omterm{prop:g4:plants-without-seeds:damp}{苔藓}、\omterm{prop:g4:plants-without-seeds:damp}{蕨类}、潮虫、蟾蜍）会减少，
而喜阳的\omterm{def:g6:exploring-our-environment:components}{生物}（耐旱的草、蚱蜢，或许还有
蜥蜴）则会迁入。

\textbf{11.} 因为重新调查与第一次只能有一点不同——
即树篱的缺失。若更换地点、
仪器或季节，就会带来各自的差异，
那么变化就无法归因于树篱了：这是
景观尺度上的公平检验原则。

\textbf{12.} 每只潮虫都会四处游走、转向，直到身体
感到舒适为止，因此这一群最终会聚集在潮湿
之处——\omterm{def:g1:animals-around-us:animal}{动物}的位置记录了\omterm{def:g6:exploring-our-environment:conditions}{条件}，
与\omterm{def:g6:exploring-our-environment:conditions}{温度}计的读数一样可靠。没有谁在驱赶它们；每只
个体自身的舒适感完成了绘图工作，这正是
为什么居民可以被当作测量数据来读取。
\end{solution}
